% Box-Muller transform: two uniforms on the unit square mapped to an
% isotropic Gaussian cloud via the polar construction of section 13.3.
\begin{figure}[htbp]
\centering
\begin{tikzpicture}[every node/.style={font=\small}]

% ---- Left: uniform square (source) ----
\begin{scope}[shift={(0, 0)}]
  \draw[thick, draw=engblack] (0, 0) rectangle (3, 3);
  % a scattering of uniform points (hand-placed, illustrative)
  \foreach \x/\y in {0.4/2.6, 1.1/0.7, 2.2/2.1, 0.8/1.9, 2.7/0.5,
    1.6/2.8, 0.3/1.2, 2.4/1.4, 1.9/0.4, 0.6/2.2, 2.9/2.7, 1.3/1.5,
    2.0/2.5, 0.9/0.9, 1.7/1.1, 2.6/1.9} {
    \fill[engblue] (\x, \y) circle (1.3pt);
  }
  \node[anchor=north] at (1.5, -0.2) {$(U_{1}, U_{2})$ uniform};
  \node[anchor=south, font=\scriptsize] at (1.5, 3.05)
    {source: independent uniforms};
\end{scope}

% ---- Middle: arrow with the transform ----
\begin{scope}[shift={(3.6, 1.5)}]
  \draw[-{Stealth}, very thick, draw=engblack!75] (0, 0) -- (1.6, 0);
  \node[anchor=south, font=\scriptsize, align=center] at (0.8, 0.1)
    {$R=\sqrt{-2\ln U_{1}}$\\[-1pt] $\Theta=2\pi U_{2}$};
  \node[anchor=north, font=\scriptsize] at (0.8, -0.1)
    {Jacobian $r \cdot \tfrac1r = 1$};
\end{scope}

% ---- Right: isotropic Gaussian cloud (target) ----
\begin{scope}[shift={(6.0, -1.5)}]
  \draw[->, thick] (-3, 0) -- (3, 0) node[right, font=\scriptsize] {$X$};
  \draw[->, thick] (0, -3) -- (0, 3) node[above, font=\scriptsize] {$Y$};
  % concentric circles = density contours
  \draw[draw=engred!50] (0, 0) circle (1.0);
  \draw[draw=engred!35] (0, 0) circle (1.8);
  % a scattering of Gaussian points clustered near origin
  \foreach \x/\y in {0.3/0.2, -0.5/0.4, 0.8/-0.3, -0.2/-0.7, 0.1/1.1,
    -1.0/0.1, 0.6/0.9, -0.4/-0.3, 1.2/0.5, -0.7/-1.0, 0.2/-0.5,
    -0.9/0.8, 0.5/0.1, -0.1/0.6, 1.5/-0.4, -1.3/-0.6, 0.0/0.0,
    0.7/-1.2, -0.6/1.3, 0.9/0.2} {
    \fill[engblue] (\x, \y) circle (1.3pt);
  }
  \node[anchor=north] at (0, -3.1) {$(X, Y)$ standard normal};
\end{scope}

\end{tikzpicture}
\caption[The Box-Muller transform]{The Box-Muller construction of
section~\ref{sec:vol02:ch13:spaces} (developed in the change-of-variables
subsection). Two independent uniforms on the unit square (left) map to a
pair of independent standard normals (right) through the polar
construction in the middle. The radius squared is exponential and the
angle is uniform; the factor $r$ from the polar area element and the
factor $1/r$ from the Cartesian Jacobian cancel, leaving the isotropic
Gaussian whose circular contours are drawn at one and two standard
deviations. Points are hand-placed for illustration; the companion
generator produces the true sample. Source: standard construction,
\cite[ch.~4]{text:ross-probability}.}
\label{fig:vol02:ch13:box-muller}
\end{figure}
